\scanpage{93}
\noindent\hfill\handwritten{\unsure{27|22}.8.1967.}\\

\noindent Cher Dieudonné,\\

Tu as bien fait de tiquer sur la généralisation que je suggérai de EGA I 10.10.5. Je pense que l'analogue nonnoethérien de l'équivalence de a) et c) doit sortir sans mal, mais en tous cas ces conditions ne son\unsure{t} pas impliquées par b). P\unsure{q|g|\deleted{q}o|\deleted{g}o}ur le voir, soit $A=B\handwritten{[[}T\handwritten{]]}$ , et définissons
\[
M_n\handwritten{\unsure{(}}=A_n/J_{\deleted{\unsure{p|y}}n}\ ,\quad\text{où}\quad A_n=B\handwritten{[}T\handwritten{]}/(T^{n+1})\ ,\quad J_n=(X_o,X_1T,..,X_nT^n\handwritten{\unsure{)}}
\]
en choisissant l'anneau $B$ et les $X_n\handwritten{\in}B$ de telle façon que les générateurs écrits de $J_n$ soient en nombre minimal. IL suffit pour ceci de prendre pour $B$ un anneau de polynômes
\[
B=k\handwritten{[}(X_i)_{1\handwritten{\geq}0}\handwritten{]}\ ,
\]
on vérifie en effet que les générateurs écrits donnent alors dans $J_n/TJ_n$ des éléments linéairement indépendants sur $\deleted{B/T}\ A_o=A/TA=B$ . En effet, si les $a_i\handwritten{\in}B$ $(0\handwritten{\leq}i\handwritten{\leq}n)$ sont tels que
\[
a_oX_o+\ldots+\deleted{\unsure{t}}a_nX_nT^n=F_oX_oT+\ldots+F_nX_nT^{n+1}\ ,
\]
où les $F_n$ sont dans $B\handwritten{[}T\handwritten{]}=A\handwritten{[}T,(X_i)\handwritten{]}$ , on voit pour tout $i$, $0\handwritten{\leq}i\handwritten{\leq}n$, divisant par l'idéal engendré par les $X_j$ avec $j\handwritten{\ne}i$ , qu'on a une relatio\unsure{n}
\[
a_iX_iT^i=F'_iX_iT^{i+1}\quad\text{dans}\quad k\handwritten{[}X_i,T\handwritten{]},
\]
où $F'_i$ est l'image de $F_i$ . Comme $X_i$ et $T$ sont des éléments réguliers de cet anneau, on en conclut que
\[
a_i=F'_i\,T\quad\text{dans}\quad k\handwritten{[}X_i,T\handwritten{]}\ ,
\]
d'où évidemment $a_i=0$ , cqfd.

Je suggère donc de rédiger la nouvelle proposition 10.10.5 en disant qu'on a deux conditions équivalentes \deleted{\unsure{en|et}} impliquant une troisième condition, et que cette dernière équivaut aux deux autres dans le cas noethérien. Il serait peut-être utile de donner aussi le contre-exemple qui précède.
% archive sheet number: 92
